\documentclass[aspectratio=169,10pt]{beamer}

\usetheme{Madrid}
\usecolortheme{default}
\setbeamertemplate{navigation symbols}{}
\setbeamertemplate{footline}[frame number]

\usepackage{amsmath,amssymb,amsthm,booktabs,mathtools}
\usepackage{graphicx}
\usepackage{hyperref}
\usepackage{tikz}
\usetikzlibrary{arrows.meta,positioning,decorations.pathreplacing}
\usepackage{xeCJK}
\setCJKmainfont{Songti SC}
\setCJKsansfont{Heiti SC}
\setCJKmonofont{Heiti SC}
\XeTeXlinebreaklocale "zh"
\XeTeXlinebreakskip = 0pt plus 1pt
% 生成 ActualText，否则字体子集化后 PDF 的中文文本层无法复制
\XeTeXgenerateactualtext=1

\definecolor{positive}{HTML}{0173B2}
\definecolor{negative}{HTML}{DE8F05}
\definecolor{emphasis}{HTML}{029E73}
\definecolor{neutral}{gray}{0.55}
\definecolor{cbPurple}{HTML}{CC78BC}
\newcommand{\pos}[1]{\textcolor{positive}{#1}}
\newcommand{\con}[1]{\textcolor{negative}{#1}}
\newcommand{\HL}[1]{\textcolor{emphasis}{#1}}

\title[Awesome GFlowNets]{Awesome GFlowNets}
\subtitle{从流匹配到可认证训练：213 篇文献的结构化盘点}
\author{awesome\_Gflow 调研流水线}
\institute{检索截止 2026-08-25}
\date{2026 年 9 月}

\begin{document}

\begin{frame}
  \titlepage
\end{frame}

\begin{frame}{本次调研的交付物}
  \begin{columns}[T]
    \column{0.48\textwidth}
      \begin{block}{四类产物}
        \begin{itemize}
          \item \textbf{213} 篇论文分类目录（TSV + markdown）
          \item \textbf{35} 篇核心论文中文深度解读
          \item \textbf{34} 篇保版式中文翻译 PDF（32 全文 + 2 篇长材料章节节选）
          \item \textbf{4} 份趋势与失效边界报告
        \end{itemize}
      \end{block}
    \column{0.48\textwidth}
      \begin{block}{证据纪律}
        \begin{itemize}
          \item 解读的论断标注论文章节或公式号
          \item 个人判断显式标注为「我的判断」
          \item 趋势报告每条给来源链接与日期
          \item 检索不到写「未检索到」，不猜
        \end{itemize}
      \end{block}
  \end{columns}
  \vspace{6pt}
  技术栈：SuperTranslate 保版式引擎 $+$ 自建 vLLM（Qwen2.5-32B-AWQ，A100 80G，16K 上下文）。
\end{frame}

\begin{frame}{GFlowNet 解决的问题：采样而非最大化}
  训练前向策略 $P_F$，使终止对象 $x$ 的采样概率满足
  \[
    P_\top(x) \;\propto\; R(x),
  \]
  其中 $R(x) > 0$ 是终端奖励，$P_\top$ 是策略诱导的终止分布。

  \vspace{8pt}
  \begin{center}
  \begin{tikzpicture}[
    box/.style={draw, rounded corners, minimum width=2.1cm, minimum height=0.75cm,
                font=\footnotesize, fill=positive!10},
    arr/.style={-{Stealth}, thick}
  ]
    \node[box] (A) {$s_0$ 空对象};
    \node[box, right=1.1cm of A] (B) {$s_1$ 加片段};
    \node[box, right=1.1cm of B] (C) {$\cdots$};
    \node[box, right=1.1cm of C] (D) {$x$，奖励 $R(x)$};
    \draw[arr] (A) -- node[above=3pt, font=\scriptsize] {$P_F$} (B);
    \draw[arr] (B) -- (C);
    \draw[arr] (C) -- (D);
  \end{tikzpicture}
  \end{center}

  \vspace{4pt}
  同一个 $x$ 可由多条轨迹到达（合流 DAG），这是与自回归模型的核心分野：需要流守恒约束才能保证边缘分布正确。

  \vspace{4pt}
  \pos{为什么重要}：药物发现要的是结构多样的候选池，单个「最优分子」在湿实验里很可能失败。\con{奖励最大化方法天然坍缩到单模态}，这是结构性差异而非调参问题。
\end{frame}

\begin{frame}{训练目标谱系：信用传播距离与方差的折中}
  \vspace{4pt}
  \begin{center}
  \begin{tabular}{llll}
    \toprule
    目标 & 约束粒度 & 解决什么 & 代价 \\
    \midrule
    Flow Matching (T01) & 单状态入出流 & 最初的正确性保证 & 长轨迹信用传播慢 \\
    Detailed Balance (T02) & 单边 & 不需枚举父节点 & 比值梯度方差大 \\
    \textbf{Trajectory Balance} (T03) & 整条轨迹 & 一步传到起点 & 长轨迹乘积、需估 $Z$ \\
    SubTB (T05) & 子轨迹（$\lambda$ 加权） & DB 与 TB 间插值 & 终止漂移、有效性退化 \\
    $f$-TB / 散度族 (T49/T17) & 轨迹级 $f$-散度 & 纳入更大家族 & 散度选择成新超参 \\
    \bottomrule
  \end{tabular}
  \end{center}

  \vspace{8pt}
  TB 的轨迹级恒等式（$\tau = (s_0 \to \cdots \to s_n = x)$）：
  \[
    Z \prod_{t=1}^{n} P_F(s_t \mid s_{t-1}) \;=\; R(x) \prod_{t=1}^{n} P_B(s_{t-1} \mid s_t)
  \]
  $Z$ 是可学习的配分函数标量，$P_B$ 是后向策略。两侧取对数作差平方即 TB loss。
\end{frame}

\begin{frame}{趋势一：从「平衡族」走向「证书族」与「博弈族」}
  2026 年的新论文不再只提新平衡式，而是附带可验证保证或引入对抗机制。

  \vspace{6pt}
  \begin{columns}[T]
    \column{0.31\textwidth}
      \textbf{\pos{证书族}}：Stable GFlowNets（arXiv 2605.01729）打通 TB loss 与全局 TV 误差，给出 loss 到 TV 的上界与有限样本证书；反例证明小 TV 误差不蕴含有界 loss，后期 loss 尖峰可能良性。
    \column{0.31\textwidth}
      \textbf{\pos{博弈族}}：DTB / ACE（arXiv 2602.17827）训练探索策略专攻当前网络低估的区域，与主网络构成双人博弈并刻画均衡，多样高奖励区域发现率超 curiosity 基线。
    \column{0.31\textwidth}
      \textbf{\pos{LLM 专用}}：RapTB $+$ SubM（arXiv 2603.00454）针对前缀坍缩与长度偏置，把子轨迹监督锚在根节点、吸收后缀回传奖励，做到 32B 规模。
  \end{columns}

  \vspace{8pt}
  \begin{alertblock}{读法}
    单纯提出新平衡恒等式的窗口正在关闭；增量空间在「目标 $+$ 可验证保证」与「目标 $+$ 探索机制」。
  \end{alertblock}
\end{frame}

\begin{frame}{趋势二：loss 到分布误差的理论链条已闭合}
  \vspace{2pt}
  \begin{center}
  \begin{tikzpicture}[
    box/.style={draw, rounded corners, minimum width=2.4cm, minimum height=0.7cm,
                font=\scriptsize, fill=emphasis!10, align=center},
    arr/.style={-{Stealth}, thick}
  ]
    \node[box] (A) {TB loss $\downarrow$};
    \node[box, right=0.9cm of A] (B) {KL 上界\\{\tiny N001}};
    \node[box, right=0.9cm of B] (C) {TV 上界 $+$ 证书\\{\tiny 2605.01729}};
    \node[box, right=0.9cm of C] (D) {可认证训练};
    \draw[arr] (A) -- (B); \draw[arr] (B) -- (C); \draw[arr] (C) -- (D);
  \end{tikzpicture}
  \end{center}

  \vspace{4pt}
  \begin{center}
  \begin{tabular}{lll}
    \toprule
    结果 & 内容 & 来源 \\
    \midrule
    收敛速率分层 & FM 达 $\mathcal{O}(1/\sqrt{T})$，DB 仅 $\mathcal{O}(1/T^{1/3})$ & arXiv 2505.02035 \\
    样本复杂度 & $\mathcal{O}\!\left(|S| L \log(|S|/\delta)/\epsilon^2\right)$ & 同上 \\
    隐式正则 & FM 偏最大熵，DB 偏 KL，TB 偏路径长度效率 & 同上 \\
    奖励噪声鲁棒性 & 按 $R_{\min}^{-4}$ 缩放 & 同上 \\
    泛化界 & 轨迹越长、目标越尖则越难 & ICLR 2025 \\
    \bottomrule
  \end{tabular}
  \end{center}

  \vspace{4pt}
  \con{尚未闭合}：以上速率均依赖 tabular 或可实现性假设，函数逼近下的收敛性未检索到任何结果。
\end{frame}

\begin{frame}{趋势三：与 GRPO 系的硬对比}
  Scaffold-conditioned SMILES 优化（数据来自 RapTB 论文 Table 1）：
  \vspace{4pt}
  \begin{center}
  \begin{tabular}{lrrrrr}
    \toprule
    方法 & 有效性 & 奖励 & 熵 & 指纹多样性 & 平均长度 \\
    \midrule
    PPO   & 1.000 & 0.604 & $\approx 0$ & --- & --- \\
    GRPO  & 0.997 & 0.661 & 0.98 & --- & 10.0 \\
    TB    & 0.998 & 0.717 & 2.503 & 0.807 & 3.065 \\
    SubTB & 0.328 & 0.755 & 2.127 & 0.836 & 8.354 \\
    \textbf{RapTB + SubM} & 0.988 & \textbf{0.844} & \textbf{2.726} & \textbf{0.898} & 7.435 \\
    \bottomrule
  \end{tabular}
  \end{center}

  \vspace{6pt}
  肽生成压力测试（bioRxiv 2026-01）：GFlowNet 二肽采样均匀度高 \pos{$5.4\times$}、奖励方差低 \pos{$1.9\times$}、重复序列少 \pos{$3.9\times$}。去掉奖励中的熵门控后，GRPO-D 的 1000 个样本 \con{$100\%$ 含同一三肽模式}，GFlowNet 多样性保持 $0.937$。
\end{frame}

\begin{frame}{护城河判断}
  \begin{exampleblock}{核心结论}
    GFlowNet 的护城河不是峰值奖励，而是\HL{对奖励函数设计缺陷的鲁棒性}。
  \end{exampleblock}

  \vspace{6pt}
  \begin{itemize}
    \item 上表中 RapTB+SubM 的奖励分数（0.844）也高于 GRPO（0.661）——差距不止在多样性。
    \item GRPO 系需要精调的多样性惩罚与熵门控才不坍缩，而这些超参在真实药物发现流水线里通常调不准。
    \item GRPO 综述（arXiv 2603.06623）把 diversity preservation 单列为研究方向（DiverseGRPO 用 Vendi Score 提升 13\%--18\%），从反面印证这是结构性弱点而非实现问题。
  \end{itemize}
\end{frame}

\begin{frame}{相对位置：不是升降，是性质变化}
  \begin{columns}[T]
    \column{0.48\textwidth}
      \textbf{\con{下降的维度}}
      \begin{itemize}
        \item 纯采样性能：物理场景 SOTA 已被 iDEM、力引导漂移、FP++ 占据，它们用能量或力的结构，GFN 的离散组合优势用不上。
        \item 理论独立性：熵正则 RL（T14）、Markov chain、OT（O08）、连续时间路径测度四个视角都能把 GFlowNet 写成特例。
      \end{itemize}
    \column{0.48\textwidth}
      \textbf{\pos{上升的维度}}
      \begin{itemize}
        \item 目标采用度：扩散采样器方向把 TB/VarGrad $+$ 回放 $+$ 局部搜索当作标准配置；torchgfn 把 off-policy 训练写进官方指南。
        \item 应用护城河：抗奖励设计缺陷在分子与肽两个独立实验都有量化证据。
        \item FlowRL 作为 recipe 合并进 verl（PR \#3924，2025-11-06），首次进入 2 万 star 级生产框架。
      \end{itemize}
  \end{columns}
\end{frame}

\begin{frame}{失效边界：什么时候 GFlowNet 不 work}
  213 篇论文里只有 1 篇的简介提到失效——证据被埋在实验章节与 limitations 里。
  \vspace{4pt}
  \begin{center}
  \begin{tabular}{llp{5.4cm}l}
    \toprule
    失效模式 & 触发条件 & 定量表现 & 依据 \\
    \midrule
    奖励峰值过尖 & $\sigma \le 0.5$ & 只找 1--2 个模态，MCMC 找 11--16 个，成本还高 $500\times$ & N081 \\
    截断致流泄漏 & 稀疏奖励 & 强制终止态违反流守恒，偏向最长轨迹；稠密奖励下无影响 & N113 \\
    系统性欠拟合 & 高熵初始化 & $\log p$ 对 $\log R$ 回归斜率仅 0.58（理想 1.0） & T07 \\
    LLM 前缀坍缩 & 终端奖励 $+$ 回放偏置 & 早期 token 熵骤降，长度系统性偏离 & T54 \\
    loss 与误差脱钩 & 任意 & 小 TV 误差不蕴含有界 loss & T51 \\
    \bottomrule
  \end{tabular}
  \end{center}
\end{frame}

\begin{frame}{元失效：四个流行评测指标都不测分布正确性}
  \begin{alertblock}{被否证的指标（T32）}
    Spearman 相关（$\log p_\theta = c\log\tilde\pi$ 时恒为 1.0，但只有 $c=1$ 才匹配目标）、平均未归一化目标、Shen's accuracy、发现模态数与 top-$k$ 奖励——全都会给「在模态上堆积过多质量」的错误分布高分。
  \end{alertblock}
  \vspace{2pt}
  唯一被验证可用的是 \textbf{FCS}（Flow Consistency in Sub-graphs）：与 TV 距离的 Spearman 相关 \pos{0.99}（集合生成）/ \pos{0.90}（序列生成），计算量少约 \pos{3 个数量级}，并有 PAC 界。

  \vspace{4pt}
  用 FCS 复检的负面结果：LED- 与 FL-GFlowNets 在「发现高价值状态」上显著优于标准 TB，但 \con{采样分布是错的}——用模态数评测会把它们判为优胜。

  \vspace{4pt}
  \textbf{操作规则}：声称「分布匹配更好」的实验，不报 FCS（或可枚举时的精确 TV）就没有说服力。
\end{frame}

\begin{frame}{工具生态的冷静一面（GitHub API 当日实测）}
  \begin{columns}[T]
    \column{0.48\textwidth}
      \textbf{\con{停滞信号}}
      \begin{itemize}
        \item Recursion/Valence 2025-06-27 开源 synflownet-boltz，此后零提交。
        \item 主库 recursionpharma/gflownet 一年仅 2 次提交（Dockerfile/tox 杂务）；公司 2025-06 裁员 20\% 后口径转向 foundation model。
        \item TRL 与 OpenRLHF 对 GFlowNet 的 issue/PR 搜索命中均为 0（对照：GRPO in TRL $=$ 1650）。
      \end{itemize}
    \column{0.48\textwidth}
      \textbf{\pos{健康信号}}
      \begin{itemize}
        \item torchgfn 是唯一全指标健康的库：v2.4.1（2026-04-05），2026-04 起提交 $\ge 100$ 次。
        \item verl 集成 FlowRL 是孤例但确实成立。
        \item gfnx（JAX，2025-11）提出标准化评测意图，但 PyPI 停在 0.0.1。
      \end{itemize}
  \end{columns}
  \vspace{4pt}
  FlowRL 自身已被 GFlowRL（arXiv 2607.13394）判定在 MoE 235B 配置下配分函数网络不收敛，改用 in-batch Monte Carlo 估 $\log Z$。
\end{frame}

\begin{frame}{最值得投入：GFlowNet $\times$ 最优传输}
  \begin{exampleblock}{已建立的结果（O08）}
    最小流非无环 GFlowNet 学习问题可等价写成线性规划；固定初始边流分布后，即变为以图上距离为代价的 Kantorovich OT 问题（也等价于图上 Beckmann 问题的离散形式）。最优前向策略采样的正是最优路径，诱导最优耦合。
  \end{exampleblock}

  \vspace{4pt}
  对偶方向：DDSBM（ICLR 2025）把 Iterative Markovian Fitting 扩展到离散空间解 Schrödinger 桥，其动力学等价于以\textbf{图编辑距离}为代价的熵正则 OT。两条线在 2026 年是收敛而非替代：SB/OT 给目标与代价结构，GFlowNet 给离散组合空间上的可扩展参数化与 off-policy 训练。

  \vspace{4pt}
  \textbf{三个具体切口}
  \begin{itemize}
    \item 把 O08 的 LP 等价性从「固定初始边流」放宽到一般情形，检验最优耦合是否仍可由 GFlowNet 策略表示。
    \item 用 GFlowNet 求解精确解不可行规模的图 OT，与 Sinkhorn 类方法做代价-时间曲线对比。
    \item 把 DDSBM 的「最小图变换」性质引入 GFlowNet 分子优化：保持 $P_\top(x) \propto R(x)$ 的同时约束编辑距离。
  \end{itemize}
\end{frame}

\begin{frame}{仓库结构与工程化}
  \begin{columns}[T]
    \column{0.48\textwidth}
      \begin{block}{内容组织（按 awesome-ml4co 规范）}
        \begin{itemize}
          \item \texttt{papers/}：213 篇 TSV 目录 $+$ 核心清单
          \item \texttt{notes/}：8 节固定模板的中文解读
          \item \texttt{pdfs/en}, \texttt{pdfs/zh}：原文与译文对照
          \item \texttt{insights/}：四份趋势与失效边界报告
          \item \texttt{surveys/}：理论指南与 OT 潜力分析
        \end{itemize}
      \end{block}
    \column{0.48\textwidth}
      \begin{block}{全流程脚本化，可增量重跑}
        \begin{itemize}
          \item \texttt{download\_core\_pdfs.py}：arXiv 检索落盘，断点续传
          \item \texttt{translate\_batch.py}：并行翻译，已完成自动跳过
          \item \texttt{generate\_notes.py}：PDF 全文经 vLLM 生成初稿
          \item \texttt{build\_readme.py}：解析目录生成分区 README
        \end{itemize}
      \end{block}
  \end{columns}
  \vspace{4pt}
  解读模板固定 8 节：一句话、问题与动机、方法核心、理论结果、实验与证据、局限与批判、与谁对话、对后续研究的启示。
\end{frame}

\begin{frame}{结论与下一步}
  \begin{columns}[T]
    \column{0.31\textwidth}
      \textbf{判断一}：理论已成熟。loss 到 KL 到 TV 的链条闭合，速率分层与样本复杂度都有结果，新目标要过审必须附带可验证保证。
    \column{0.31\textwidth}
      \textbf{判断二}：护城河是鲁棒性而非峰值。应主打奖励函数难以精调的真实流水线场景。
    \column{0.31\textwidth}
      \textbf{判断三}：生态薄，工程红利未被吃掉。只有 torchgfn 健康、只有 verl 集成。
  \end{columns}

  \vspace{10pt}
  \begin{alertblock}{三个行动项}
    \begin{enumerate}
      \item 攻函数逼近下的收敛性——当前唯一明确的理论空位。
      \item 沿 O08 推进 GFN $\times$ OT：放宽 LP 等价性条件 $+$ 大规模图 OT 求解器对比。
      \item 补齐生态：面向 LLM 后训练的 GFlowNet 训练库与标准评测套件。
    \end{enumerate}
  \end{alertblock}
\end{frame}

\begin{frame}{References}
  \small
  \begin{thebibliography}{9}
    \bibitem{tb} Malkin et al. Trajectory Balance: Improved Credit Assignment in GFlowNets. NeurIPS 2022.
    \bibitem{found} Bengio et al. GFlowNet Foundations. JMLR 2023.
    \bibitem{stable} Stable GFlowNets with Probabilistic Guarantees. arXiv:2605.01729, 2026.
    \bibitem{raptb} Rooted Absorbed Prefix Trajectory Balance with Submodular Replay. arXiv:2603.00454, 2026.
    \bibitem{secrets} Secrets of GFlowNets' Learning Behavior: A Theoretical Study. arXiv:2505.02035, 2025.
    \bibitem{ot} Your GFlowNet Secretly Learns an Optimal Transport Plan. ICML 2026 SPIGM Workshop.
    \bibitem{offpolicy} Sendera et al. Improved off-policy training of diffusion samplers. NeurIPS 2024.
  \end{thebibliography}
\end{frame}

\begin{frame}
  \begin{center}
    \Large 完整目录、逐篇解读与中英文 PDF\\[8pt]
    \href{https://github.com/asimfish/awesome_Gflow}{\texttt{github.com/asimfish/awesome\_Gflow}}
    \vspace{14pt}

    \normalsize
    生成方式与已知局限见仓库 \texttt{DELIVERY.md}
  \end{center}
\end{frame}

\end{document}
